% Volume III: Attention and Publicity
% Part VI. Privacy as Boundary and Capacity
% Chapter 34: Privacy and Future Reinterpretation
% Spine: proof-spines/ch034-spine.md

\chapter{Privacy and Future Reinterpretation}
\label{ch:ch034}

Let \(M_t(R)\) denote the interpretation of a record \(R\) at time \(t\). Then
\[
M_t(R) \neq M_{t+\Delta}(R)
\]
\ProposedTheorem\ even when the record itself is unchanged. Privacy therefore protects not only present secrets, but future vulnerability to unknown classificatory regimes. A trace disclosed under one social ontology may acquire a different meaning later, under new norms, institutions, or analytic tools.

The chapter adds a temporal limit to the prior chapters on mosaics, contextual transfer, model capture, and underdescription. Because records can be reinterpreted and recombined indefinitely, privacy sometimes requires forgetting, expiration, and friction against perpetual reuse. Otherwise every preserved appearance remains permanently available for renewed suspicion. The same protection against future reinterpretation becomes concealment when expiration or forgetting rules erase records that are still needed to reconstruct or challenge exercises of power.
